\bookchapter{第三十一章：短視頻拍攝與內容創作入門}

\chapterintro{2026年，短視頻已非“年輕人之玩具”矣，它變成矣所有人獲取信息與表達自我之基礎方式。無論汝乃欲記錄生活、分享知識，抑作點副業，手機即汝最好之創作工具——不須買相機、不須學Premiere，一部手機遂能開始。此一章幫汝自零起步。}

\booksection{一、拍前先欲清楚：汝要拍賜誰觀？}

甚多人拍視頻之第一個錯誤即“拿起手機遂拍”，拍完不知道賜誰觀。於開拍之前，先回答三個問題：

\begin{enumerate}
  \item \textbf{吾之觀衆乃誰？}——拍賜孩子觀、拍賜同齡人觀、拍賜長輩觀，內容風格完全不同。
  \item \textbf{吾能提供何價值？}——使人笑（娛樂）、使人學（知識）、使人饞（美食）、使人暖（情感）。至少佔一樣。
  \item \textbf{吾能於哪個平臺持續更新？}——抖音/快手（泛娛樂）、小紅書（生活方式）、B站（深度內容）、視頻號（社交圈）。選一個主陣地深耕，勿一上來遂全平臺分發。
\end{enumerate}

\booksection{二、手機攝基設——開拍前必作之5事}

\begin{enumerate}
  \item \textbf{分辨率設1080P，幀率設30fps}：此乃短視頻平臺之標準。4K文件太大上傳慢，60fps於社交平臺上無優勢反而文件更大。
  \item \textbf{擦乾淨鏡頭}：第一條亦乃最重要之一條。拍之前用衣服擦一下鏡頭，畫面立刻清晰一個檔次。
  \item \textbf{鎖定對焦與曝光}：於取景框長按畫面主體2秒，將出現"AE/AF鎖定"提示。此樣拍之時未能蓋因移動導致畫面忽明忽暗。
  \item \textbf{橫拍抑豎拍？}——發抖音/速手/視頻號/小紅書\textbf{豎着拍}（9:16），發B站/YouTube\textbf{橫着拍}（16:9）。欲清楚平臺再決定方向。
  \item \textbf{開網格線}：設置 > 相機 > 網格線。三分法構圖之基礎，幫汝快速尋到畫面重心。
\end{enumerate}

\booksection{三、運鏡五招——使畫面“動”起來}

新手拍視頻最大之問題乃畫面太“靜”——全程站着不動拍。學會下面五個基本運鏡，畫面立刻有專業感：

\begin{itemize}
  \item \textbf{推（Push）}：手機慢慢向爲拍對象靠近。製造"靠近"之緊張感或親密感。
  \item \textbf{拉（Pull）}：手機慢慢遠離爲拍對象。自細節拉至全景，交代環境。
  \item \textbf{搖（Pan）}：手機處不動，左右轉動。適拍風景、展示空間。
  \item \textbf{移（Track）}：手機平行移動，與着爲拍對象走。須手持穩定，最好配合手機穩定器。
  \item \textbf{與（Follow）}：手機與於爲拍對象後面或前面走。Vlog最常用之運鏡。
\end{itemize}

\textbf{穩乃第一要務}：手抖之畫面看起來甚不專業。雙手持機比單手穩；走路時膝蓋微屈，用“忍者步”減少上下顛簸；若有預算，買一個手機穩定器（大疆OM系列約500-800元），效果質變。

\booksection{四、剪輯入門——剪映遂夠矣}

剪映（CapCut之國內版）乃目前最適合新手之手機剪輯工具，免費、功能全、模板多。以下乃基礎工作流：

\begin{enumerate}
  \item \textbf{導入素材}：打開剪映 > 開始創作 > 選擇拍好之視頻片段導入。
  \item \textbf{裁}：拖動時間軸，以每個片段開頭結尾之廢片剪掉。原則：每一段唯保留有信息量之部分，寧短勿長。
  \item \textbf{加轉場}：兩段視頻之間點"轉場"按鈕，選一個簡單之轉場效果。新手建議用"疊化"或"黑場"，勿用花哨之3D轉場。
  \item \textbf{添字幕}：剪映有"識別字幕"能，一鍵自動生字幕。準率約90\%，手校對一遍即可。
  \item \textbf{配音樂}：剪映音樂庫裏有大量版權免費之音樂。選一首節奏與內容匹配之BGM，音量調至視頻原聲之1/3左右。
  \item \textbf{導出}：1080P導出，幀率30fps，碼率選"薦"即可。
\end{enumerate}

\booksection{五、封面之要性——70\%之人僅唯封面決觀不觀}

短視頻之封面即“門面”，決定別人要勿點進來。幾個作封面之要點：

\begin{itemize}
  \item \textbf{封面要有一句話}：用大字寫明視頻之核心內容，比如"3分鐘學會手機修圖"。字體要大、顏色要亮、與背景有對比。
  \item \textbf{封面與內容要一致}：勿作"標題黨"——封面曰一件事，點進去乃另一回事，觀衆將立刻劃走，算法將降權。
  \item \textbf{人物表情要生動}：若封面有人，選一個表情生動或誇張之瞬間。人臉比文字更能吸引點擊。
  \item \textbf{剪映作封面}：於剪映裏選一幀最好之畫面，加文字疊加，導出爲圖片作爲封面。
\end{itemize}

\booksection{六、內容創作之三條心法}

\begin{itemize}
  \item \textbf{前3秒定生死}：短視頻平臺，觀衆3秒內決定拂去抑繼續觀。前3秒須賜出"鉤子"——一個問題、一個懸念、一個意外畫面。
  \item \textbf{一條視頻唯講一件事}：新手最容易犯之錯乃一條視頻塞太多內容。一條視頻唯解決一個問題，講透比講全重要。
  \item \textbf{堅持較天賦要}：算法推薦須"養號"。前20條視頻幾未能有播放量，此乃正常之。堅持每週更新2-3條，使算法認識汝之內容方向，流量將慢慢起來。
\end{itemize}

\begin{tcbpractice}{動手：拍一條30秒之短視頻}
\begin{enumerate}
  \item 選一個汝擅長之小知識（比如"何以將手機字體調大"）。
  \item 以手機豎屏攝3-5段素材，每段5-10秒。
  \item 用剪映綴、裁、添字幕、配BGM。
  \item 作一張帶大字之封面圖。
  \item 導出後先發予朋友觀，收集反饋再發布。
\end{enumerate}

\end{tcbpractice}

\begin{tcbsummary}{本章小結}
\begin{itemize}
  \item 拍之前欲清楚：觀衆乃誰、提供何價值、於哪個平臺發
  \item 開拍5事：1080P/30fps、擦鏡頭、鎖對焦曝光、確橫豎、開網格線
  \item 運鏡五招：推、拉、搖、移、與——穩乃第一要務
  \item 剪輯用剪映：裁轉場字幕BGM導出，六步走
  \item 前3秒定生死，一條視頻唯講一件事，堅持比天賦重要
\end{itemize}

\end{tcbsummary}
